\documentclass{article}

\usepackage{fancyhdr}
\usepackage{amsmath}
\allowdisplaybreaks
\usepackage{amsthm}
\usepackage{amsfonts}
\usepackage{tikz}
\usepackage[plain]{algorithm}
\usepackage{algpseudocode}
\usepackage{enumitem}
\usepackage{tikz}
\usepackage{pgfplots}
\pgfplotsset{compat=1.18}
\usepgfplotslibrary{groupplots}
\usepackage{comment}

\usetikzlibrary{automata,positioning}

%
% Basic Document Settings
%

\topmargin=-0.45in
\evensidemargin=0in
\oddsidemargin=0in
\textwidth=6.5in
\textheight=9.0in
\headsep=0.25in

\linespread{1.1}

\pagestyle{fancy}
\lhead{\hmwkAuthorName}
\chead{} % empty
\rhead{\hmwkHeaderTitle}
\cfoot{\thepage}

\renewcommand\headrulewidth{0.4pt}
\renewcommand\footrulewidth{0pt}

\setlength\parindent{0pt}

%
% Homework Details
%

\newcommand{\hmwkTitle}{Homework 9}
\newcommand{\hmwkDueDate}{Monday, March 9, 2026}
\newcommand{\hmwkClass}{ICMA 216 Calculus IIIA}
\newcommand{\hmwkClassTime}{Section 1}
\newcommand{\hmwkClassInstructor}{Ruth J. Skulkhu}
\newcommand{\hmwkAuthorName}{\textbf{Jiraroj Wiruchpongsanon (6781617)}}
\newcommand{\hmwkHeaderTitle}{ICMA 216 Calculus IIIA: Homework 9}

%
% Title Page
%

\title{
    \vspace{2in}
    \textmd{\textbf{\hmwkClass:\ \hmwkTitle}}\\
    \normalsize\vspace{0.1in}\small{Due\ on\ \hmwkDueDate}\\
    \vspace{0.1in}\large{\textit{\hmwkClassInstructor\ \hmwkClassTime}}
    \vspace{3in}
}

\author{\hmwkAuthorName}
\date{}

\renewcommand{\part}[1]{\textbf{\large Part \Alph{partCounter}}\stepcounter{partCounter}\\}

%
% Helper Commands
%

\newcounter{partCounter}
\newcommand{\solution}{\textbf{\large Solution}}

\newtheorem*{theorem}{Theorem}

\theoremstyle{remark}
\newtheorem*{remark}{Remark}

\begin{document}

\maketitle
\newpage

\section*{Problem 1 (13.7 Tangent Planes)}
Find the equations of the tangent plane at the point $(-2, 1, -3)$ to the ellipsoid
\[
    \frac{x^2}{4} + y^2 + \frac{z^2}{9} = 3.
\]

\medskip
\solution

first check that $(-2,1,-3)$ is on the surface of the ellipsoid:
\[
    \frac{(-2)^2}{4} + 1^2 + \frac{(-3)^2}{9} = 1 + 1 + 1 = 3
\]

now let
\[
    F(x,y,z) = \frac{x^2}{4} + y^2 + \frac{z^2}{9} - 3 = 0
\]

then a normal vector to the tangent plane is given by $\nabla F$, where
\[
    F_x = \frac{x}{2}, \qquad F_y = 2y, \qquad F_z = \frac{2z}{9}
\]

at the point $(-2,1,-3)$,
\[
    \nabla F(-2,1,-3) = \left\langle -1, 2, -\frac{2}{3} \right\rangle
\]

we can multiply by $-3$ to get a simpler normal vector
\[
    \vec n = \langle 3, -6, 2 \rangle
\]

therefore, the tangent plane through $(-2,1,-3)$ is
\[
    3(x+2) - 6(y-1) + 2(z+3) = 0
\]

which gives
\[
    3x - 6y + 2z + 18 = 0
\]

so the equation of the tangent plane is
\[
    \boxed{3x - 6y + 2z + 18 = 0}
\]

\vspace{1em}
\hrulefill

\section*{Problem 2 (13.7 Tangent Planes)}
Find an equation for the tangent plane and parametric equations for the normal line to the surface
\[
    z = x^2y
\]
at the point $(2,1,4)$. What is the acute angle that the tangent plane at the point $(2,1,4)$ makes with the $xy$-plane?

\medskip
\solution

rewrite the surface as
\[
    F(x,y,z) = x^2y - z = 0
\]

then the tangent plane at $(x_0,y_0,z_0)$ is given by
\[
    F_x(x_0,y_0,z_0)(x-x_0) + F_y(x_0,y_0,z_0)(y-y_0) + F_z(x_0,y_0,z_0)(z-z_0) = 0
\]

where
\[
    F_x = 2xy, \qquad F_y = x^2, \qquad F_z = -1
\]

at $(2,1,4)$,
\[
    F_x(2,1,4) = 2(2)(1) = 4
\]
\[
    F_y(2,1,4) = (2)^2 = 4
\]
\[
    F_z(2,1,4) = -1
\]

therefore, the tangent plane is
\[
    4(x-2) + 4(y-1) - (z-4) = 0
\]

so
\[
    \boxed{4x + 4y - z = 8}
\]


    \medskip
    \begin{center}
    \rule{0.3\linewidth}{0.2pt}
    \end{center}
    \medskip


this plane has normal vector
\[
    \vec n = \langle 4,4,-1 \rangle
\]

thus the normal line through $(2,1,4)$ has parametric equations
\[
    x = 2 + 4t, \qquad y = 1 + 4t, \qquad z = 4 - t
\]

so the normal line is
\[
    \boxed{x = 2 + 4t, \qquad y = 1 + 4t, \qquad z = 4 - t}
\]


    \medskip
    \begin{center}
    \rule{0.3\linewidth}{0.2pt}
    \end{center}
    \medskip


for the acute angle between the tangent plane and the $xy$-plane, we use the angle between their normal vectors.

the normal vector of the $xy$-plane is
\[
    \vec n_1 = \langle 0,0,1 \rangle
\]

so
\[
    \cos \theta = \frac{|\vec n \cdot \vec n_1|}{\|\vec n\|\,\|\vec n_1\|}
    = \frac{|\langle 4,4,-1 \rangle \cdot \langle 0,0,1 \rangle|}{\sqrt{4^2+4^2+(-1)^2}\cdot 1}
    = \frac{1}{\sqrt{33}}
\]

therefore,
\[
    \boxed{\theta = \arccos\left(\frac{1}{\sqrt{33}}\right)}
\]

\vspace{1em}
\hrulefill

\section*{Problem 3 (13.8 Maximum and Minimum Values)}
Find all the critical points of
\[
    f(x,y) = x^3y + 12x^2 - 8y
\]

\medskip
\solution

find the first partial derivatives:
\[
    f_x(x,y) = 3x^2y + 24x
\]
\[
    f_y(x,y) = x^3 - 8
\]

for critical points, solve
\[
    f_x = 0 \qquad \text{and} \qquad f_y = 0
\]

start with $f_y = 0$:
\[
    x^3 - 8 = 0
\]
\[
    x^3 = 8
\]
\[
    x = 2
\]

substitute this into $f_x = 0$:
\[
    3(2)^2y + 24(2) = 0
\]
\[
    12y + 48 = 0
\]
\[
    y = -4
\]

therefore, the only critical point is
\[
    \boxed{(2,-4)}
\]

\end{document}
